\chapter{Excessive Yawning}

\section*{Definition}
Excessive yawning is yawning that occurs more often than expected, beyond what is normal for tiredness or boredom. While occasional yawning is normal, persistent or frequent yawning without obvious cause may signal an underlying medical issue.

\section*{What Happens in Your Body}
Yawning is a reflex involving a deep inhalation, jaw opening, and slow exhalation. It is thought to help cool the brain, increase oxygen levels, and regulate alertness. When you yawn excessively, it may be your body's response to low oxygen, changes in blood pressure, or stimulation of the vagus nerve (the nerve that connects your brain to your chest and abdomen). Certain conditions that affect the brain's temperature regulation or blood flow can also trigger frequent yawning.

\section*{Causes}
\begin{itemize}
  \item Sleep deprivation, insomnia, or poor sleep quality (most common)
  \item Anxiety or stress
  \item Medication side effects: Some antidepressants (SSRIs), opioids, and stimulants can cause excessive yawning
  \item Vasovagal reactions: A drop in heart rate and blood pressure that can cause yawning before fainting
  \item Neurological conditions: Multiple sclerosis, migraine with aura, epilepsy (particularly temporal lobe seizures), Parkinson's disease, brain tumors
  \item Heart conditions: Some people yawn excessively before a heart attack or with aortic dissection
  \item Sleep disorders: Narcolepsy, sleep apnea, and other conditions that disrupt deep sleep
  \item Opioid withdrawal
  \item Boredom or fatigue (normal, not excessive)
\end{itemize}

\section*{When to See a Doctor}
See your doctor if:
\begin{itemize}
  \item Yawning is frequent, persistent, and unexplained by lack of sleep
  \item Yawning is accompanied by chest pain, shortness of breath, or lightheadedness
  \item You have other neurological symptoms such as vision changes, weakness, or numbness
  \item You experience sudden, excessive yawning along with any signs of a heart attack (chest pressure, arm pain, sweating)
  \item Yawning started after you began a new medication
  \item You also have excessive daytime sleepiness or sudden sleep attacks
\end{itemize}

\section*{Related Symptoms}
\begin{itemize}
  \item Fatigue
  \item Daytime sleepiness
  \item Difficulty concentrating
  \item Shortness of breath
  \item Dizziness or lightheadedness
  \item Anxiety
\end{itemize}
\textbf{Important:} Excessive yawning is rarely an emergency on its own, but if it appears with chest pain, shortness of breath, or sudden neurological symptoms, seek emergency care.

\section*{Self-Care / Home Management}
\begin{itemize}
  \item Improve sleep hygiene — aim for 7-9 hours of quality sleep per night with a consistent bedtime.
  \item Practice relaxation techniques such as deep breathing or meditation to reduce anxiety-driven yawning.
  \item Try cooling down — splash cold water on your face or step outside for fresh air.
  \item Take breaks during long periods of sedentary or repetitive activity.
  \item Review medications with your doctor if yawning started after a new prescription.
\end{itemize}

\section*{How Your Doctor Will Evaluate You}
\begin{itemize}
  \item Your doctor will ask about your sleep patterns, stress levels, medications, and whether yawning happens at specific times or with certain activities.
  \item They will check your vital signs, including heart rate and blood pressure.
  \item If a neurological cause is suspected, they may perform a basic neurological exam (reflexes, coordination, eye movements).
  \item A sleep study (polysomnography) may be ordered if sleep apnea or narcolepsy is suspected.
  \item MRI of the brain or EEG may be needed if there are other neurological symptoms or if seizures are suspected.
\end{itemize}

\section*{Treatment Options}
\begin{itemize}
  \item Treatment depends entirely on the underlying cause — excessive yawning itself is not treated directly.
  \item If caused by a medication, your doctor may adjust the dose or switch to an alternative.
  \item If caused by sleep apnea, treatment with CPAP may resolve the yawning.
  \item If related to anxiety or stress, therapy, relaxation techniques, or anti-anxiety medication may help.
  \item Neurological or cardiac causes require specific treatment directed at the underlying condition.
  \item In most cases, addressing the root cause stops the excessive yawning.
\end{itemize}


\section*{References}
\begin{thebibliography}{99}
\bibitem{walusinski2016}
Walusinski O. History of yawning. In: \textit{Mystery of Yawning in Physiology and Disease}. Karger; 2010.
\bibitem{guggisberg2010}
Guggisberg AG, Mathis J, Schnider A, Hess CW. Why do we yawn? \textit{Neurosci Biobehav Rev}. 2010;34(8):1267--1276.
\bibitem{provan1995}
Provan JW, Engel K. Yawning as a symptom of neurological disease. \textit{Neurology}. 1995;45(12):2220--2222.
\bibitem{dakin2014}
Dakin CJ, Cooper PE, Khot S. Yawning and its clinical significance. \textit{Br J Hosp Med}. 2014;75(7):404--407.
\bibitem{gallup2007}
Gallup AC, Gallup GG Jr. Yawning as a brain cooling mechanism: nasal breathing and forehead cooling diminish the incidence of contagious yawning. \textit{Evol Psychol}. 2007;5(1):92--101.
\bibitem{daquin2014}
Daquin G, Micallef J, Blin O. Yawning. \textit{Sleep Med Rev}. 2001;5(4):299--312.
\end{thebibliography}
\clearpage
